\documentclass[11pt]{scrartcl}

\usepackage[stixtwo]{fontsetup}
\usepackage{amsmath}

\begin{document}
$x^2 - 9 = (x - 3)(x + 3)$ és $x^2 - x - 6 = (x - 3)(x + 2)$, így
$x \ne 3, -2$ esetén \(f(x) = \frac{x + 3}{x + 2}\).
\begin{enumerate}
	\item Igazolni kell, hogy minden $\varepsilon > 0$ számhoz létezik
	      $\delta(\varepsilon) > 0$, hogy $0 < |x + 1| < \delta(\varepsilon)$
	      esetén $|f(x) - 2| < \varepsilon$. Legyen $\varepsilon > 0$.
	      \[
		      \left| \frac{x + 3}{x + 2} - 2 \right|
		      = \frac{|x + 3 - 2x - 4|}{|x + 2|}
		      = \frac{|x + 1|}{|x + 2|}
	      \]
	      Ha $|x + 1| < \frac{1}{2}$, akkor $-\frac{3}{2} < x < -\frac{1}{2}$, így
	      $|x + 2| > \frac{1}{2}$, tehát
	      \[
		      \frac{|x + 1|}{|x + 2|} < 2 |x + 1| < \varepsilon,
	      \]
	      ha $|x + 1| < \frac{\varepsilon}{2}$. Így
	      \(
	      \delta(\varepsilon) = \min\left( \frac{1}{2}, \frac{\varepsilon}{2} \right)
	      \).

	\item \[
		      \lim_{x \to -3} f(x)
		      = \lim_{x \to -3} \frac{x + 3}{x + 2}
		      = \frac{0}{-1}
		      = 0
	      \]

	\item \[
		      \lim_{x \to +\infty} f(x)
		      = \lim_{x \to +\infty} \frac{1 + \frac{3}{x}}{1 + \frac{2}{x}}
		      = 1
	      \]
\end{enumerate}
\end{document}
